\message{ !name(Hamilton2024.tex)}%%%%%%%%%%%%%%%%%%%%%%%%%%%%%%%%%%%%%%%%%%%%%%%%%%%%%%%%%%%%%%%%%%%%%%%%%%%% 
% AGUJournalTemplate.tex: this template file is for articles formatted with LaTeX
%
% This file includes commands and instructions
% given in the order necessary to produce a final output that will
% satisfy AGU requirements, including customized APA reference formatting.
%
% You may copy this file and give it your
% article name, and enter your text.
%
% guidelines and troubleshooting are here: 

%% To submit your paper:
\documentclass[draft]{agujournal2019}
\usepackage{url} %this package should fix any errors with URLs in refs.
\usepackage{lineno}
\usepackage[inline]{trackchanges} %for better track changes. finalnew option will compile document with changes incorporated.
\usepackage{wasysym}
\usepackage{soul}

\linenumbers
%%%%%%%
% As of 2018 we recommend use of the TrackChanges package to mark revisions.
% The trackchanges package adds five new LaTeX commands:
%
%  \note[editor]{The note}
%  \annote[editor]{Text to annotate}{The note}
%  \add[editor]{Text to add}
%  \remove[editor]{Text to remove}
%  \change[editor]{Text to remove}{Text to add}
%
% complete documentation is here: http://trackchanges.sourceforge.net/
%%%%%%%

\draftfalse

%% Enter journal name below.
%% Choose from this list of Journals:
%
% JGR: Atmospheres
% JGR: Biogeosciences
% JGR: Earth Surface
% JGR: Oceans
% JGR: Planets
% JGR: Solid Earth
% JGR: Space Physics
% Global Biogeochemical Cycles
% Geophysical Research Letters
% Paleoceanography and Paleoclimatology
% Radio Science
% Reviews of Geophysics
% Tectonics
% Space Weather
% Water Resources Research
% Geochemistry, Geophysics, Geosystems
% Journal of Advances in Modeling Earth Systems (JAMES)
% Earth's Future
% Earth and Space Science
% Geohealth
%
% ie, \journalname{Water Resources Research}

\journalname{Paleoceanography and Paleoclimatology}


\begin{document}

\message{ !name(Hamilton2024.tex) !offset(-3) }


%%%%%%%%%%%%%%%%%%%%%%%%%%%%%%%%%%%%%%%%%%%%%%%
%  TITLE
%
% (A title should be specific, informative, and brief. Use
% abbreviations only if they are defined in the abstract. Titles that
% start with general keywords then specific terms are optimized in
% searches)
%
%%%%%%%%%%%%%%%%%%%%%%%%%%%%%%%%%%%%%%%%%%%%%%%

% Example: \title{This is a test title}

\title{Interpreting benthic Common Era foraminiferal records under the assumption of modern-day circulation}

%%%%%%%%%%%%%%%%%%%%%%%%%%%%%%%%%%%%%%%%%%%%%%%
%
%  AUTHORS AND AFFILIATIONS
%
%%%%%%%%%%%%%%%%%%%%%%%%%%%%%%%%%%%%%%%%%%%%%%%

% Authors are individuals who have significantly contributed to the
% research and preparation of the article. Group authors are allowed, if
% each author in the group is separately identified in an appendix.)

% List authors by first name or initial followed by last name and
% separated by commas. Use \affil{} to number affiliations, and
% \thanks{} for author notes.
% Additional author notes should be indicated with \thanks{} (for
% example, for current addresses).

% Example: \authors{A. B. Author\affil{1}\thanks{Current address, Antartica}, B. C. Author\affil{2,3}, and D. E.
% Author\affil{3,4}\thanks{Also funded by Monsanto.}}

\authors{=list all authors here=}


% \affiliation{1}{First Affiliation}
% \affiliation{2}{Second Affiliation}
% \affiliation{3}{Third Affiliation}
% \affiliation{4}{Fourth Affiliation}

\affiliation{=number=}{=Affiliation Address=}
%(repeat as many times as is necessary)


% Corresponding author mailing address and e-mail address:

% (include name and email addresses of the corresponding author.  More
% than one corresponding author is allowed in this LaTeX file and for
% publication; but only one corresponding author is allowed in our
% editorial system.)

% Example: \correspondingauthor{First and Last Name}{email@address.edu}

\correspondingauthor{=name=}{=email address=}



%%%%%%%%%%%%%%%%%%%%%%%%%%%%%%%%%%%%%%%%%%%%%%%
% KEY POINTS
%%%%%%%%%%%%%%%%%%%%%%%%%%%%%%%%%%%%%%%%%%%%%%%
%  List up to three key points (at least one is required)
%  Key Points summarize the main points and conclusions of the article
%  Each must be 140 characters or fewer with no special characters or punctuation and must be complete sentences

% Example:
% \begin{keypoints}
% \item	List up to three key points (at least one is required)
% \item	Key Points summarize the main points and conclusions of the article
% \item	Each must be 140 characters or fewer with no special characters or punctuation and must be complete sentences
% \end{keypoints}

\begin{keypoints}
\item We present an inversion method that incorporates age model uncertainty to infer far-field properties from benthic sediment core records (137/140 characters)
 \item The reconstructed North Atlantic Little Ice Age is characterized by a 0.11-1.2K/kyr cooling (1$\sigma$ confidence interval) from 1200-1900CE (140/140 characters) 
\item A long-term cooling in the Nordic Seas is reconstructed, and is found to be approximately two times greater than that of the subpolar gyre (138/140 characters)
\end{keypoints}

%%%%%%%%%%%%%%%%%%%%%%%%%%%%%%%%%%%%%%%%%%%%%%%
%
%  ABSTRACT and PLAIN LANGUAGE SUMMARY
%
% A good Abstract will begin with a short description of the problem
% being addressed, briefly describe the new data or analyses, then
% briefly states the main conclusion(s) and how they are supported and
% uncertainties.

% The Plain Language Summary should be written for a broad audience,
% including journalists and the science-interested public, that will not have 
% a background in your field.
%
% A Plain Language Summary is required in GRL, JGR: Planets, JGR: Biogeosciences,
% JGR: Oceans, G-Cubed, Reviews of Geophysics, and JAMES.
% see http://sharingscience.agu.org/creating-plain-language-summary/)
%
%%%%%%%%%%%%%%%%%%%%%%%%%%%%%%%%%%%%%%%%%%%%%%%

%% \begin{abstract} starts the second page

\begin{abstract}
  Paleoceanographic reconstructions of the past 2,000 years indicate a large-scale cooling in the latter half of the last millennium, but significant conflict on magnitude and timing exist. This conflict is partially due to the inference of surface conditions from sparse, noisy, near-surface marine proxies that can exhibit seasonal biases.
Here, we develop a novel methodology for inferring far-field changes from benthic sediment core records via an ocean circulation model. Benthic sediment core records are not subject to high-energy near-surface variability or seasonal biases and record much more stable signals than equivalent planktonic records. 
We present a possible surface temperature and $\delta^{18}\mathrm{O}_{\mathrm{seawater}}$ history that fits a dataset of 11 high-resolution benthic $\delta^{18}\mathrm{O}_{\mathrm{calcite}}$ records.
The estimate presented here is within error of instrumental records, and is characterized by a Little Ice Age sea surface temperature cooling of 0.67 $\pm$ 0.56 K/kyr, which is similar to, but at the high end, of previous independent estimates obtained from Common Era marine surface records.
Our reconstruction is characterized by a larger magnitude millenial-scale cooling in the Nordic Seas (0.9 $\pm$ 2.7 K) than in the subpolar gyre (0.5 $\pm$ 2.2K), indicative of an Arctic amplification of Little Ice Age cooling.
 

\end{abstract}

\section*{Plain Language Summary}
Enter your Plain Language Summary here or delete this section.
Here are instructions on writing a Plain Language Summary: 
https://www.agu.org/Share-and-Advocate/Share/Community/Plain-language-summary


%%%%%%%%%%%%%%%%%%%%%%%%%%%%%%%%%%%%%%%%%%%%%%%
%
%  BODY TEXT
%
%%%%%%%%%%%%%%%%%%%%%%%%%%%%%%%%%%%%%%%%%%%%%%%

%%% Suggested section heads:
\section{Introduction}
The past 2,000 years, referred to as the Common Era (CE), is a critical period for understanding the response of the Earth's climate to perturbations on long timescales. Climatic perturbations during the CE include decreasing solar irradiance due to the precession of the equinoxes and stratospheric sulfate injection from volcanic eruptions. During the latter half of the last millennium, a climatic anomaly referred to as the Little Ice Age has been identified, characterized by anomalously cool sea surface temperatures \cite{ocean2k} as well as surface air temperatures \cite[]{Consortium--2013:Continental}. Understanding the spatiotemporal extent of the Little Ice Age is a critical step in constraining the response of the climate system to various natural perturbations on timescales longer than that of the instrumental record. Commonly, the Little Ice Age is thought to have been triggered by either decreasing solar irradiance or volcanic activity, and reinforced via oceanic and atmospheric mechanisms. However, the precise reinforcement mechanisms have yet to be conclusively identified due to the sparse, noisy, and biased nature of paleoclimate proxy records. 

Modeling studies have suggested that an increase of Arctic Ocean sea-ice as a direct response to decadally-paced volcanic eruptions beginning in 1257-1258CE may have lead to a freshening of the subpolar North Atlantic that weakened convection \cite{Zhong-Miller-2010:Centennial, miller12, Lehner2013, schleussner2013}. Weakened convection has been associated with a slowdown of the Atlantic Meridional Overturning Circulation (AMOC), although some modeling studies have suggested that an AMOC slowdown is not necessary to explain the full magnitude of LIA cooling \cite[]{Moreno-Chamarro-Zanchettin-2016:Little}. In addition, paleoceanographic proxy records have not identified large-scale circulation anomalies in the pre-industrial CE \cite{Thornalley-Oppo-2018:Anomalously, moffasanchez2015, Rahmstorf-Box-2015:Exceptional}. Dynamics related to a weakened subpolar gyre \cite{Moffa-Sanchez-Born-2014:Solar, Moreno-Chamarro-Zanchettin-2016:Little} have also been suggested to be a key reinforcement mechanism, and have been supported through deep-sea coral records \cite{Copard-Colin-2012:Late}, bivalve resiliency records \cite{Arellano-Nava2022}, and planktonic foraminiferal records \cite[]{Moffa-Sanchez-Born-2014:Solar, moffasanchez2017, Spooner-Thornalley-2020:Exceptional}, although exact timing and spatial extent of subpolar gyre variability still remain unclear. 

% Studies are inconclusive on the timing of the LIA, with start date estimates ranging from 1250-1380. Few constraints on the magnitude of the cooling of the LIA exist, with published estimates of . Warming from the 1870s to 2000s has been estimated to be 0.59 $\pm$ 0.12 $^{\circ}$C \cite[]{Roemmich-John-2012:Challenger}.
%Cunningham: 0.02$^{\circ}$C/century 
One of the primary limitations of our understanding of LIA dynamics comes from the point-wise nature of paleoceanographic proxy records. Inferring behavior of a large-scale oceanographic feature such as the subpolar gyre typically requires spatially dense and high temporal resolution measurements, which can be difficult to attain for paleo-timescales. Estimates of LIA sea surface temperature (SST) cooling rates as well as timing differ amongst reconstructions, from an estimate of 0.2$^{\circ}$C/kyr from 1200-1940CE in the northeast North Atlantic \cite[]{Cunningham-Austin-2013:Reconstruction}, to 0.54$^{\circ}$C/kyr from 800-1900CE globally \cite{ocean2k}. These reconstructions all rely upon surface records, which experience a high degree of noise and seasonal effects. When trying to extract long-term climate signals, benthic records may be a better option, as they are not subject to high-frequency surface variability or seasonal variability. 

In this paper, we develop an inverse approach to inferring far-field properties from benthic proxy records through the usage of a data-informed circulation model. We apply a modern-day circulation model under the assumption that there were no large-scale circulation changes over the CE. Testing whether or not simultaneous sediment core records can be well-fit under the assumption of steady-state modern-day circulation has been previously used as a test for the necessity of circulation change \cite[]{Burke-Marchal-2011:Application, Gebbie-2012:Tracer, Amrhein-Gebbie-2015:Inferring}. The circulation model used here is a transient inversion of modern-day tracer observations \cite[]{Gebbie-Huybers-2012:mean} previously used to constrain CE deep Pacific and Atlantic cooling \cite[]{gebbie19}. This prior application estimated a temperature history that could explain the observed warming from the 1870s HMS Challenger expedition to modern hydrography. It has also been used to explain the lag between Atlantic and Pacific warming in deglacial $\delta^{18}\mathrm{O}$ records \cite[]{Gebbie-2012:Tracer}. Here, we apply this model to CE benthic sediment core $\delta^{18}\mathrm{O}$ records for the first time. 

We will first show the ability of this methodology to fit a novel dataset of 11 high-resolution benthic records of $\delta^{18}\mathrm{O}_{\mathrm{calcite}}$ from along the Reykjanes Ridge in Iceland \cite[]{Lu-Oppo-2023:Surface}. Successfully fitting this data implies that possible surface temperature and $\delta^{18}\mathrm{O}_{\mathrm{seawater}}$ histories exist that can explain the sediment core data without necessitating changes in circulation.
Next, we will compare our reconstruction to instrumental records of Labrador Sea Water temperature and salinity, and demonstrate that we reconstruct similar magnitudes of temperature and salinity variability over the past 100 years.
Our reconstruction can also be compared to other paleoceanographic estimates, and we show that the estimated LIA North Atlantic cooling trend is comparable to the global SST trend of 0.54$^{\circ}$C/kyr identified in the Ocean2k dataset.
Finally, we will discuss the spatiotemporal evolution of the LIA in the North Atlantic, and quantify the regional magnitudes and timings. 
In our solution, the Nordic Seas are shown to be the first region to begin cooling, at $\sim$1200CE, followed by the subpolar gyre in $\sim$1400CE. The timing of the subpolar gyre cooling is similar to other independent records \cite[]{https://doi.org/10.1002/2013PA002523}. The subpolar gyre is shown to cool by 0.37 $\pm$ 2.1K, while the Nordic Seas cool by 0.8 $\pm$ 2.7 K. This finding is consistent with the Arctic amplification of the LIA previously identified in the Community Earth System Model's Last Millenium Ensemble (CESM-LME) \citeA[]{ottobliesner16}. Arctic amplification was previously suggested to be exhibited in this dataset by larger magnitude $\delta^{18}\mathrm{O}_{\mathrm{calcite}}$ variations in benthic records than in planktonic records \cite{Lu-Oppo-2023:Surface}. Arctic amplification has previously been identified in paleoclimate records from the Holocene Thermal Maximum, Last Glacial Maximum, and mid-Pliocene \cite[]{Miller-Alley-2010:Arctic}, but has not been previously conclusively identified in Common Era records. 

\section{Data}
\subsection{Sediment Core Records from Lu et al. (2023) \label{sec:ludata}}
The inversion presented here is applied to a set of 11 high-resolution benthic foraminiferal $\delta^{18}\mathrm{O}_{\mathrm{calcite}}$ sediment core records (Figure \ref{fig:recon}) previously presented in \citeA[]{Lu-Oppo-2023:Surface}. These records were collected from along the Reykjanes Ridge south of Iceland, and form a depth transect from 1100-2300m. Previously, these records were used to demonstrate that the deep ocean LIA transition, as indicated by the benthic records, occurred at 1350 $\pm$ 50 CE, while the surface ocean LIA transition, as indicated by planktonic records, occurred from 1100-1700CE. In addition, a larger magnitude LIA was observed in the benthic records than in the planktonic records, indicative of Arctic amplification of LIA temperature changes, transferred to the sediment core sites by an active AMOC \cite[]{Lu-Oppo-2023:Surface}.

\citeA{Lu-Oppo-2023:Surface} used two benthic species of foraminifera to measure $\delta^{18}\mathrm{O}_{\mathrm{calcite}}$. The $\delta^{18}\mathrm{O}_{\mathrm{calcite}}$ measurements were obtained from Uvigerina peregrina (U. peregrina) tests for the three shallowest cores (MC28A, MC26A, and MC25A), and from Cibicidoides wuellerstorfi (C. wuellerstorfi) tests for the other eight sediment cores.

The BACON software package in R was used to generate a Bayesian accumulation history for each sediment core based on 3-8 radiocarbon samples and an assigned coretop age (see \citeA{Lu-Oppo-2023:Surface} for more details). Each sediment core has an age model that consists of 1000 realizations of age with respect to depth. 

  
\subsection{Total Matrix Intercomparison (TMI), steady-state and transient versions}
In order to relate signals at the sediment core sites to far-field ocean regions, an ocean circulation model must be applied. Here, we apply the results of \citeA{gebbie2010}, in which a box inversion technique is used to compute the pathways connecting interior ocean grid boxes to the surface. The global ocean is decomposed into a $4^{\circ} \times 4^{\circ}$ grid with 33 vertical levels, and the pathways that connect each interior grid box to the surface are computed by inverting six conservative tracers. This results in $\mathbf{A}$, a matrix that maps between global surface tracer distributions and the corresponding global interior tracer distribution assuming the ocean is at a modern-day steady-state circulation. The matrix $\mathbf{A}$ consists of mass fractions from each surface grid cell to every interior cell. In \citeA{Gebbie-Huybers-2012:mean}, this result is extended by incorporating radiocarbon observations, resulting in a discrete time tracer transport operator $\mathbf{L}$ which is similar to $\mathbf{A}$ but represents mass fluxes instead of fractions. 

\subsection{Coupled Climate Model Output}
[NOTE: I feel like this doesn't really need to be introduced, maybe I could just distribute this into Section \ref{sec:Cuu}]
SST fields from historical (1850-2015CE) realization `r11i1p1f1` of the Community Earth System Model, version 2 (CESM2) \cite[]{cesm2} were used to constrain the relative variance of surface spatial modes. 

An estimate of the relationship between temperature and $\delta^{18}\mathrm{O}_{\mathrm{seawater}}$ was obtained using output from the iCESM1.2 coupled climate model \cite[]{Brady-Stevenson-2019:Connected}. The \texttt{TEMP} and \texttt{R18O} variables from the large ensemble pre-industrial fully coupled control run were observed at the grid cell closest to the mean sediment core position, and the linear least squares relationship of the two timeseries was found to be 10$^{\circ}$C/$\permil$. 

\section{Methods}
\subsection{Overview \label{sec:overview}}
This inverse problem is characterized by inferring some time-varying far-field properties $\mathbf{u}$ that best explain data at the sediment core sites $\mathbf{y}$.
Both $\mathbf{y}$ and $\mathbf{u}$ are vectorized arrays, with $\mathbf{y}$ consisting of some sediment core property indexed by time $t_{y}$ and sediment core $c$, and with $\mathbf{u}$ indexed by time $t_{u}$, spatial mode $m$, and state variables known to affect whatever sediment core property is reported in $\mathbf{y}$. The connection between the sediment core sites and any region of the ocean can be computed via the discrete time transport operator $\mathbf{L}$. 

We are concerned with solving for the time-varying evolution of some collection of far-field patterns. These patterns could be oceanographic regions or features, and could be surface or subsurface. In order to reduce the number of degrees of freedom in the problem, we solve for the magnitude of surface spatial modes, obtained from a singular value decomposition of source water maps (see Section \ref{sec:spatmode}). These spatial modes can be propagated through the discrete time transport operator $\mathbf{L}$ and ``observed'' at the core sites, forming the dimensional array $\mathbf{\mathcal{M}}$ (Sections \ref{sec:spatmode} and \ref{sec:E}). $\mathbf{\mathcal{M}}$ consists of the magnitude of spatial modes propagated to the core sites at any lag time from some initial signal. Once $\mathbf{\mathcal{M}}$ is formed, an impulse response technique can be used to obtain a model matrix $\mathbf{E}$, which provides the basis of the inverse problem (Equation \ref{eq:inverse}).

\begin{equation}
\label{eq:inverse}
\mathbf{y} = \mathbf{Eu}
\end{equation}

We solve Equation \ref{eq:inverse} for the time-varying magnitudes of surface patterns $\mathbf{u}$ by using a tapered weighted least squares approach. This methodology minimizes the cost function shown in Equation \ref{eq:costfunc}, where $\mathbf{\tilde{u}}$ represents values of $\mathbf{u}$ that best minimize $J$ and $\mathbf{u_0}$ is the prior, or first guess, of $\mathbf{\tilde{u}}$. The vector $\mathbf{\tilde{y}}$ is the estimated values at the core sites according to $\mathbf{\tilde{u}}$ (e.g. $\mathbf{\tilde{y}} = \mathbf{E\tilde{u}}$). Finally, the matrices $\mathbf{C_{yy}}$ and $\mathbf{C_{u_0u_0}}$ are covariance matrices defining the expected statistics of $\mathbf{y}$ and $\mathbf{u_0}$, respectively. 

\begin{equation}  
\label{eq:costfunc}
J = \left(\mathbf{y}-\mathbf{\tilde{y}}\right)\mathbf{C_{yy}}^{-1}\left(\mathbf{y}-\mathbf{\tilde{y}}\right)^{T}+(\mathbf{\tilde{u} - u_0})\mathbf{C_{u_{0}u_{0}}}^{-1}(\mathbf{\tilde{u} - u_0})^T
\end{equation}

The formula for a tapered, weighted, least squares solution and the solution covariance $\mathbf{C_{\tilde{u}\tilde{u}}}$ is shown in Equation \ref{eq:twls}. 
 
\begin{equation}
\mathbf{\tilde{u}} = \mathbf{u_{0}} + \mathbf{C_{u_{0}u_{0}}} \mathbf{E}^T\left(\mathbf{E}\mathbf{C_{u_{0}u_{0}}}\mathbf{E}^T+\right)^{-1}\left(\mathbf{y}-\mathbf{Eu_0}\right)
\label{eq:twls}
\end{equation}
\begin{equation}
\mathbf{C_{\tilde{u}\tilde{u}}} = \mathbf{C_{u_{0}u_{0}}}\mathbf{E}^T \left(\mathbf{E}\mathbf{C_{u_{0}u_{0}}}\mathbf{E}^T\right)^{-1}\mathbf{C_{yy}} \left(\mathbf{E}\mathbf{C_{u_{0}u_{0}}}\mathbf{E}^T\right)^{-1}\mathbf{E}\mathbf{C_{u_{0}u_{0}}}
    \label{eq:twlsCuu}
    \end{equation}


In the following sections, we detail how $\mathbf{y}$ (Section \ref{sec:agemodel} and Section \ref{sec:tanddelta}), $\mathbf{C_{yy}}$ (Section \ref{sec:agemodel}), $\mathbf{u_0}$ (Section \ref{sec:spatmode}), $\mathbf{C_{u_0u_0}}$ (Section \ref{sec:Cuu}), and $\mathbf{E}$ (Section \ref{sec:E}) are formed. 

\subsection{Incorporating Age Model Uncertainty \label{sec:agemodel}}
Here, we incorporate age model uncertainty into the inverse problem by applying the BACON-generated age models \cite[]{Bacon, McKay-Emile-Geay-2021:R} from \citeA{Lu-Oppo-2023:Surface}. The age models are used to generate mean time series of $\delta^{18}\mathrm{O}_{\mathrm{calcite}}$ at evenly-spaced time intervals for $\mathbf{y}$, in addition to defining the statistics of the associated covariance matrix $\mathbf{C_{yy}}$.

Each age model realization was linearly interpolated at a common 10-year resolution. Then, for each 10-year timestep, the mean $\delta^{18}\mathrm{O}_{\mathrm{calcite}}$ was computed for any interpolated timestep that had more than 100 realizations available, populating $\mathbf{y}$. The covariance between two arbitrary timesteps $t_{1}, t_{2}$ can be computed according to Equation \ref{eq:covariance}, where $\delta^{18}\mathrm{O}_{i}(t_1)$ refers to the interpolated $\delta^{18}\mathrm{O}_{\mathrm{calcite}}$ value at $t_{1}$ for realization $i$, $<\delta^{18}\mathrm{O}(t_1)>$ is the mean value across all realizations at timestep $t_1$, and $M$ is the total number of realizations with interpolated values at $t_{1}$ and $t_{2}$. Realizations will have different start and end times due to the Bayesian age modeling process, so there are not necessarily 1000 realizations at every timestep. Finally, a laboratory uncertainty of 0.07$\permil$ was added as a diagonal element on to the covariance matrix. 

%[NOTE: This may be a good spot for a supplementary figure more completely explaining this process.]

\begin{equation}
\label{eq:covariance}
  \mathbf{C_{yy}}[t_1, t_2] = \frac{\sum_i^M \left(\delta^{18}\mathrm{O}_i(t_1) - <\delta^{18}\mathrm{O}(t_1)>
    \right) \left( \delta^{18}\mathrm{O}_i(t_2) - <\delta^{18}\mathrm{O}(t_2)> \right) }{M} 
\end{equation}

\subsection{Deconvolving $\delta^{18}\mathrm{O}_{\mathrm{seawater}}$ and temperature \label{sec:tanddelta}}

The measured $\delta^{18}\mathrm{O}_{\mathrm{calcite}}$ of a foraminiferal test is dependent upon the temperature and $\delta^{18}\mathrm{O}_{\mathrm{seawater}}$ of ambient seawater at the time of calcification. The $\delta^{18}\mathrm{O}_{\mathrm{seawater}}$ is in turn affected by equilibrium fractionation processes such that higher $\delta^{18}\mathrm{O}_{\mathrm{seawater}}$ is associated with cooler and saltier waters, while lower $\delta^{18}\mathrm{O}_{\mathrm{seawater}}$ is associated with warmer and fresher waters. Ultimately, higher values of $\delta^{18}\mathrm{O}_{\mathrm{calcite}}$ can be interpreted as being associated with either cooler temperatures and/or more saline waters.

This relationship has been quantified in coretop foraminiferal samples by \citeA{marchitto}. A linear approximation for $\delta^{18}\mathrm{O}_{\mathrm{calcite}}$ as a function of temperature and $\delta^{18}\mathrm{O}_{\mathrm{seawater}}$ was established and is shown in Equation \ref{eq:cibs}. Here, the experimentally-derived relationship for C. wuellerstorfi is shown, where $\delta^{18}\mathrm{O}_{\mathrm{calcite}}$ is in $\permil$ Pee Dee Belemnite (PDB) and $\delta^{18}\mathrm{O}_{\mathrm{seawater}}$ is in $\permil$ Standard Mean Ocean Water (SMOW). While C. wuellerstorfi calcifies at equilibrium, U. peregrina is isotopically heavier by 0.47$\permil$ \cite[]{marchitto}. \citeA{marchitto} recommend removing the 0.47$\permil$ offset from U. peregrina records, and applying Equation \ref{eq:cibs}. 

\begin{equation}
\label{eq:cibs}
\delta^{18}\mathrm{O}_{\mathrm{calcite}} = (-0.224 \pm 0.002)\frac{\permil}{K} T + (3.53 \pm 0.02)\permil + (\delta^{18}\mathrm{O}_{\mathrm{seawater}} - 0.27\permil)
\end{equation}

% \begin{equation}
%   \delta^{18}\mathrm{O}_{\mathrm{calcite}} = (-0.231 \pm 0.004)\frac{\permil}{K} T + (4.03 \pm 0.03)\permil + \delta^{18}\mathrm{O}_{\mathrm{seawater}} 
% \label{eq:uvi}
% \end{equation}

In order to assemble the vector $\mathbf{y}$, the values reported from both C. wuellerstorfi and U. peregrina records must be put on a common scale. In addition, the values used to define the prior $\mathbf{u_0}$ must be on the same scale. The most straightforward way to define the prior would be based upon modern-day global surface spatial distributions of temperature and $\delta^{18}\mathrm{O}_{\mathrm{seawater}}$. For this to be possible, the first guess propagated to the coresites, $\mathbf{Eu_0}$, should be near or within the range of reported values in $\mathbf{y}$.

We can use Equation \ref{eq:cibs} to estimate effective modern-day $\delta^{18}\mathrm{O}_{\mathrm{calcite}}$ at the core sites from the global gridded estimate of $\delta^{18}\mathrm{O}_{\mathrm{seawater}}$ and temperature from the inversion of \citeA[]{gebbie2010}. We find that the total recorded range of $\delta^{18}\mathrm{O}_{\mathrm{calcite}}$ values for each sediment core are variably offset from the computed effective $\delta^{18}\mathrm{O}_{\mathrm{calcite}}$ at the core site. One important factor to recognize is that the average radiocarbon date uncertainty ($\sim$200years) is near the temporal extent of rapid industrial oceanic changes. Records may capture different temporal extents of the industrial era, and the offset from modern-day values will be highly sensitive to how much of the industrial era warming is captured. For this reason, we choose to remove the 1850-1980 mean from each sediment core record prior to forming $\mathbf{y}$. We can then establish a first guess that oceanic conditions are equivalent to the 1850-1980 mean. Since the mean 1850-1980 $\delta^{18}\mathrm{O}_{\mathrm{calcite}}$ value is 0, the first guess will accordingly be a vector of zeros. 

\begin{figure}[htbp]
\centerline{\includegraphics[width = 0.75\textwidth]{/home/brynn/Code/OPTinv.jl/plots/offsets}}
\caption[]{\label{fig:offset} Rainbow dots, with size decreasing as time increases, show mean $\delta^{18}\mathrm{O}_{\mathrm{calcite}}$ for each sediment core, with 0.47$\permil$ equilibration offset removed from U. peregrina records. Black crosses show effective $\delta^{18}\mathrm{O}_{\mathrm{calcite}}$ computed from the steady-state inversion of temperature and $\delta^{18}\mathrm{O}_{\mathrm{seawater}}$ from \citeA{gebbie2010} and observed at the sediment core sites. Computation of effective $\delta^{18}\mathrm{O}_{\mathrm{calcite}}$ is according to Equation \ref{eq:cibs}.}
\end{figure}

\subsection{Computing surface spatial modes \label{sec:spatmode}}
As stated in Section \ref{sec:overview}, we want to solve our system for some far-field surface spatial pattern or region. In order to minimize the degrees of freedom in the problem, while maximizing the amount of information our solution contains, we use the steady-state TMI matrix to quantify the spatially-varying contribution of the surface ocean to the sediment core sites. A singular value decomposition approach is used to characterize the variability between core sites in terms of surface source water origination. Once surface spatial modes are identified, they are propagated through the TMI matrix and observed at the sediment core sites, quantifying the relationship between sediment core sites and the identified surface patterns. We can then solve for the magnitude of these surface spatial modes for each state variable (in this case, temperature and $\delta^{18}\mathrm{O}_{\mathrm{seawater}}$) throughout time that best explain the sediment core data. 

The relationship between each sediment core and every surface grid cell in TMI can be obtained by computing a source water inversion according to the steady-state TMI matrix $\mathbf{A}$. We solve Equation \ref{eq:swinv} for $\mathbf{s}_{\mathrm{surface}}$, which is a vectorized surface map of source water percentage to some core site. The vector $\mathbf{\delta}$ contains the sediment core location interpolated to the TMI grid system. 

\begin{equation}
\mathbf{\delta = A s_{\mathrm{surface}}}
\label{eq:swinv}
\end{equation}

Modes characterizing the variability of the surface source water maps can be identified via singular value decomposition. Each vector $\mathbf{s_i}$ can be concatenated into a matrix, which we then column-weight by a Gaussian distance matrix with a radius of 1000km. The singular value decomposition of the resultant matrix is taken, resulting in right singular vectors $\mathbf{V}$ that highlight common regions of importance in terms of source water percentage, while maximizing the orthogonality of the modes. These spatial modes form a complete basis set for the steady-state problem, but can still be applied to the transient problem, as the only difference is propagation speed, not the spatial pattern of circulation. Modes are shown in Figure \ref{fig:svdmodes}.

\begin{figure}[htbp]
\centerline{\includegraphics[width = 0.75\textwidth]{/home/brynn/Code/OPTinv.jl/plots/SVD_modes.png}}
\caption[]{\label{fig:svdmodes} Surface spatial modes obtained from the singular value decomposition of surface source water origination maps to the core sites.}
\end{figure}

\subsection{Expected statistics of surface spatial mode magnitudes \label{sec:Cuu}}
%see `optinv_decompose.jl` for details on this 
The expected statistics of the surface spatial mode magnitudes are determined using SST fields from the historical realization of CESM2. At each timestep, the CESM2 SST fields are linearly interpolated to the TMI grid and vectorized, forming a vector $\mathbf{s_i}$ where $i$ is a time index in the annually-averaged  CESM2 `historical` run (1850-2015). We can then solve for the magnitude of each surface spatial mode $\mathbf{m_i}$ from Equation \ref{eq:Cuustats}.

\begin{equation}
\label{eq:Cuustats}
\mathbf{Vm_i} = \mathbf{s_i}
\end{equation}

The standard deviation of each surface spatial mode's magnitude with respect to time is then calculated, which will then populate the diagonal elements of $\mathbf{C_{u_0u_0}}$ for temperature. The diagonal elements of $\mathbf{C_{u_0u_0}}$ that are associated with $\delta^{18}\mathrm{O}_{\mathrm{seawater}}$ will be populated by the temporal standard deviation of $\mathbf{m_i}$, divided by 10$\frac{^{\circ}C}{\permil}$, according to the expected ratio between changes in temperature and $\delta^{18}\mathrm{O}_{\mathrm{seawater}}$.

Off-diagonal elements of $\mathbf{C_{u_0u_0}}$ will be populated according to Equation \ref{eq:offdiags}, where $\rho$ is the defined correlation coefficient and $\sigma_{i}$ refers to the temporal standard deviation of mode $i$ from CESM2. 

\begin{equation}
\label{eq:offdiags}
\mathbf{C_{u_0u_0}}[t_u = t_u, m = m, \theta = \delta] = \rho \frac{\sigma_i^2}{10 ^{\circ}C / \permil}
\end{equation}

Despite non-perfect correlation between temperature and $\delta^{18}\mathrm{O}_{\mathrm{seawater}}$ in reality, we choose to set $\rho = 1$. As there are not independent constraints on temperature and $\delta^{18}\mathrm{O}_{\mathrm{seawater}}$, but linearity is enforced via the problem definition, the solution will similarly have a linear relationship between the two variables. In order for the solution to have a linear relationship at the value suggested by climate model output analysis (10$^{\circ}$/$\permil$), $\rho$ must be set equal to 1. 

\subsection{Propagating surface spatial modes \label{sec:E}}
Following the approach in \citeA[]{Gebbie-2012:Tracer} (Appendix A), we will propagate each mode $\mathbf{v_i}$ through the transient TMI matrix by solving the differential equation in Eq. \ref{eq:diffeq}, where $\mathbf{L}$ is the discrete time tracer transport operator and $\mathbf{c}$ is the magnitude of the propagated spatial mode at each TMI grid cell, interior and exterior.

\begin{equation}
\label{eq:diffeq}
\mathbf{c}(t + \Delta t) = \mathbf{Lc}(t) 
\end{equation}

In order to solve Eq. \ref{eq:diffeq}, we prescribe a spatial mode as the initial condition and use the Dirichlet boundary condition mapping matrix $\mathbf{B}$ to map this surface condition onto a matrix with the same form as $\mathbf{c}(t)$. Therefore, the initial condition will be $\mathbf{c_0} = \mathbf{Bv_i}$. We solve Eq. \ref{eq:diffeq} for $\mathbf{c}(t)$ over some large period of time, such as 1-2000 yrs at 1 year resolution with a QNDF solver, and observe the result $\mathbf{c}(t)$ at the core site locations. This is equivalent to computing the step response to a constant-in-time boundary condition, which is much more computationally efficient than computing the response to a time-varying boundary condition. The impulse response, or the response to a surface condition that persists for 1 year, and then disappears, can be computed by taking the time derivative of the step response.

We will refer to the resultant array as $\mathbf{\mathcal{M}}$, where every value will be the propagated spatial mode at the sediment core site at a particular lagged time from some arbitrary initial time. If $\mathbf{y}$ is provided at a coarser time resolution than 1 year, $\mathbf{\mathcal{M}}$ can be subsampled by summing the contribution from each year between time steps in $\mathbf{y}$.

The relationship between the sediment core records $\mathbf{y}$ and the identified surface spatial mode magnitudes throughout time $\mathbf{u}$ is shown in Equation \ref{eq:impresp}. In order to calculate a signal at the core site according to some surface mode magnitude at a particular time $t$, we must sum the contribution of each surface mode to the core site at all previous relevant lagged times $\tau$ and for each state variable $s$. The constant $\alpha$ allows for the constant -0.224K/$\permil$ from Equation \ref{eq:cibs} to be applied to the temperature component of $\mathbf{u}$. For the $\delta^{18}\mathrm{O}_{\mathrm{seawater}}$ component of $\mathbf{u}$, $\alpha = 1$. 

\begin{equation}
\label{eq:impresp}
\mathbf{y}[t_y, c] = \sum_s \sum_{\tau} \alpha[s] \mathbf{u}[t_y - \tau, s] \mathbf{\mathcal{M}}[\tau]
\end{equation}

Individual rows of $\mathbf{E}$ can be identified by making small perturbations to an arbitrary input vector $\mathbf{u}$, applying Equation \ref{eq:impresp}, and scaling the output by the perturbation. 

\subsection{Analysis techniques}
Once we arrive at a solution $\mathbf{u}$, we can visually examine the fit of the data by computing the solution reconstruction at the core sites. This is achieved by either applying Equation \ref{eq:impresp}, or by computing $\mathbf{E\tilde{u}}$. The associated covariance matrix can be computed according to Equation \ref{eq:solcov}.
\begin{equation}
\label{eq:solcov}
\mathbf{C_{\tilde{y}\tilde{y}}} = \mathbf{E C_{\tilde{u}\tilde{u}}E}^T 
\end{equation}

Output surface maps of temperature and $\delta^{18}\mathrm{O}_{\mathrm{seawater}}$ throughout time can be constructed by using the $\mathbf{V}$ matrix, as shown in Equation \ref{eq:surfrecon}.

\begin{equation}
\label{eq:surfrecon}
  \mathbf{\tilde{s}}[s = \theta](t_u) = \mathbf{V \tilde{u}}[s = \theta](t_u) 
\end{equation}

\section{Results}
For the methodology presented here, we trim all all sediment cores to a common time period shared by all sediment cores. For all 11 records in the \citeA[]{Lu-Oppo-2023:Surface} dataset, the common time period is from 1470-1980. We also present an alternative solution, comprised of sediment cores MC26A, MC22A, and MC13A, which have a common time period of

\subsection{Data Fit}
We can examine the values of $\mathbf{\tilde{u}}$ throughout time for to understand the variability of specific spatial surface modes. In Figure \ref{fig:mags}, we show the timeseries for the first three spatial modes for both solutions. From Figure \ref{fig:svdmodes}, we note that Mode 1 consists of a single spatial anomaly of consistent sign, centered over Iceland. Mode 2 is a dipole that is positive south of Iceland and negative north of Iceland, and Mode 3 is a dipole that is positive over the Labrador Sea and negative over the Irminger Sea.

With these mode definitions in mind, we can interpret the long-term decline in Mode 1 in the \texttt{old} solution as a long-term, full-basin cooling trend. Mode 2 demonstrates centennial-scale variability, with maxima near $\sim$1350CE and $\sim$1850CE. Mode 3 demonstrates variability with lower frequency than that of Mode 2, with a maxima also near $\sim$1350CE, but a minima near $\sim$1850CE. Modes are not necessarily comparable between the two tests. As $\mathbf{y}$ has different dimensionality, the degrees of freedom will be higher in the \texttt{all} case, and variability may not be distributed across the modes comparable. However, we note that there is a long-term decrease in Mode 1 in the \texttt{all} case, that reaches its most negative value in 1800CE, and then returns to a postive value. We note that Modes 2 and 3 have centennial-scale variability, with the variability in Mode 3 having higher frequency than that of Mode 2. 

Some relevant numerical features of the solution mode magnitudes in Figure \ref{fig:mags} include the largely constant behavior prior to the existence of data in $\mathbf{y}$. The observational vector $\mathbf{y}$ and the mode magnitude vector $\mathbf{u}$ can exist over different time periods because of the lagged response of $\mathbf{y}$ to $\mathbf{u}$. However, we note that at some point, $\mathbf{u}$ no longer has any data informing it, and has largely constant behavior (e.g., pre 800CE for the \texttt{old} solution). As a result of each timestep in $\mathbf{y}$ being the sum of the earlier timesteps in $\mathbf{u}$, we note that the final timestep in $\mathbf{u}$ returns to the unconstrained values. This occurs because most of the information relevant to determining $\mathbf{u}$ at the final timestep has not had long enough to propagate to depth. 

\begin{figure}[htbp]
\centerline{\includegraphics[width = 0.75\textwidth]{/home/brynn/Code/OPTinv.jl/plots/modemagsθallvold.png}}
\caption[]{\label{fig:mags} Temperature surface mode magnitudes, scaled to a z-score, for first three modes, mode 1 in purple/solid, mode 2 in green/dashed, mode 3 in orange/dotted for \textbf{(A)} estimate using three oldest cores and \textbf{(B)} estimate using all cores.}
\end{figure}

By propagating all modes through the model matrix, we can generate timeseries of the reconstructed solution at the core site, as shown in Figure \ref{fig:recon}. Initially, we note that all solutions fall within the 1$\sigma$ errorbars of the original sediment core records, implying that the data is fairly well fit. This suggests that large-scale circulation devations from the mean is not required to explain this dataset.

Here, we note that although many of the sediment core records in close proximity have disparate mean signals (i.e., MC19A and MC10A), our inverse framework cannot permit similar behavior, and instead provides a ``middle-of-the-road solution'', in which the solution reconstruction is fairly similar between neighboring core sites. We note that in the \texttt{all} solution, the reconstructed solution has smaller magnitude changes in $\delta^{18}\mathrm{O}_{\mathrm{calcite}}$ at shallower sites than at deeper sites. Additionally, we note that the reconstruction for the ``old'' solution exhibits largely long-term increases in $\delta^{18}\mathrm{O}_{\mathrm{calcite}}$, while the ``all'' solution exhibits an increase in $\delta^{18}\mathrm{O}_{\mathrm{calcite}}$ that reaches its peak $\sim$1850CE, and then decreases towards modern-day values. This could be due to the inclusion of sediment cores that exhibit significant post-industrial changes (e.g., MC21A, MC10A). 

\begin{figure}[htbp]
\centerline{\includegraphics[width = \textwidth]{/home/brynn/Code/OPTinv.jl/plots/reconsolallvold.png}}
\caption[]{\label{fig:recon} Reconstruction at core sites for all core estimate (red) and three oldest core estimate (blue). Black line shows sediment core record. }
\end{figure}

\subsection{Solution Comparison to Instrumental Record}

% \begin{figure}[htbp]
% \centerline{\includegraphics[width = \textwidth]{/home/brynn/Code/OPTinv.jl/plots/post1900allvoldδ.png}}
% \caption[]{\label{fig:post1900S} Contours: estimated MLD salinity (computed as $\delta^{18}\mathrm{O}$ divided by 0.5255) as an anomaly from 1900-1970. Scatters: MLD salinity (estimated using Holte and Talley algorithm) from EN4 CTD profiles binned by decade, and into 10$^{\circ}$ bins, computed as anomaly from climatological mean MLD salinity in EN4 objective analysis from 1900-1970.}
% \end{figure}

% \begin{figure}[htbp]
% \centerline{\includegraphics[width = \textwidth]{/home/brynn/Code/OPTinv.jl/plots/post1900allvoldθ.png}}
% \caption[]{\label{fig:post1900T} Contours: estimated MLD temperature as an anomaly from 1900-1970. Scatters: MLD temperature (estimated using Holte and Talley algorithm) from EN4 CTD profiles binned by decade, and into 10$^{\circ}$ bins, computed as anomaly from climatological mean MLD temperature in EN4 objective analysis from 1900-1970.}
% \end{figure}

We can compute time-varying surface maps of temperature and $\delta^{18}\mathrm{O}_{\mathrm{seawaeter}}$ via Equation \ref{eq:surfrecon}, which is shown in Figure \ref{fig:sst}. Although obversational timeseries on relevant timescales are rare, we average the SST map in a Labrador Sea box, and compare it to a comparable region in HadISST, as well as a digitized long-term Labrador Sea Water observational dataset 
\begin{figure}[htbp]
\centerline{\includegraphics[width = \textwidth]{/home/brynn/Code/OPTinv.jl/plots/labcompallvold.png}}
\caption[]{\label{fig:labcomp} \textbf{(A)} Reconstructed Labrador Sea (60$^{\circ}$W-30$^{\circ}$E, 55$^{\circ}$N-64$^{\circ}$N) temperature and \textbf{(B)} salinity (computed as $0.5255 \times \delta^{18}\mathrm{O}_{\mathrm{seawater}}$) over the instrumental period (1850-1970) is shown for ``old'' (blue) and ``all'' (red) solutions. Solid line shows digitized Labrador Sea record from The Lab Sea Group, 1998 publication. Dashed line shows HadiSST Labrador Sea spatial mean (60$^{\circ}$W-30$^{\circ}$E, 55$^{\circ}$N-64$^{\circ}$N). All records are shown as an anomaly from 1970.}

\end{figure}

% \begin{figure}[htbp]
% \centerline{\includegraphics[width = \textwidth]{/home/brynn/Code/OPTinv.jl/plots/labcomp_tands_allvold.png}}
% \caption[]{\label{fig:tands} Temperature and salinity diagram for the old cores solution (left, blue) and all cores solution (right, red), averaged over 60$^{\circ}$W and 30$^{\circ}$E and 55$^{\circ}$N and 64$^{\circ}$N for 1920-1970, as an anomaly from 1970. Also shown is a decadally-averaged LSW record \cite[]{lsdce}.}
% \end{figure}

\subsection{Common Era Solution}

\begin{figure}[htbp]
\centerline{\includegraphics[width = \textwidth]{/home/brynn/Code/OPTinv.jl/plots/surfacesolallvold.png}}
\caption[]{\label{fig:sst} Reconstruction of surface potential temperature, 70-year mean, SVD approach, anomaly from 1850-1970. Top row is for old cores, middle row is for all cores, third row is Last Millenium Reanalysis SSTs \cite[]{Hakim-Emile-Geay-2016:last}.}
\end{figure}


\begin{figure}[htbp]
\centerline{\includegraphics[width = 0.75\textwidth]{/home/brynn/Code/OPTinv.jl/plots/meantsallvold.png}}
\caption[]{\label{fig:meants} Mean North Atlantic temperature for ``all'' solution (red) and ``old'' (blue), Last Millenium Reanalysis (gray dash-dot), HadiSST (black dotted). Ocean2k trend of 0.54-0.38 $^{\circ}$C/kyr (from Supplemental Figure 6) is shown in shaded gray region. All records are shown as anomaly from 1850-1970. For HadiSST and both reconstructions, the mean is computed over the region (50$^{\circ}$W-20$^{\circ}$E, 50$^{\circ}$N-90$^{\circ}N$)

[This will go in discussion, but for now: Estimated LIA N. Atl. cooling is calculated by computing the linear least squares slope between the maxima and minima occurring between 1000CE and 1900CE. Old: 0.67 $\pm$ 0.56 K/kyr, All: 0.77 $\pm$ 1.8 K/kyr] } 
\end{figure}

\begin{figure}[htbp]
\centerline{\includegraphics[width = \textwidth]{/home/brynn/Code/OPTinv.jl/plots/EGWGallvold.png}}
\caption[]{\label{fig:EGWGSI} Temperature for \textbf{(A)} ``old'' and \textbf{(B)} ``all'' solutions observed for spatial averages in Nordic Seas (purple: 20$^{\circ}$W-20$^{\circ}$E, 65$^{\circ}$N-80$^{\circ}$N), the subpolar gyre (green: 60$^{\circ}$W-0$^{\circ}$, 50$^{\circ}$N-65$^{\circ}$N), the western subpolar gyre (green, dashed: 60$^{\circ}$W-30$^{\circ}$W, 50$^{\circ}$N-65$^{\circ}$N), and the eastern subpolar gyre (green, dotted: 30$^{\circ}$W-0$^{\circ}$, 50$^{\circ}$N-65$^{\circ}$N). All reconstructions are shown as anomaly from 1850-1970CE. Blue line shows digitized GRIP borehole record \cite[]{Dahl-Jensen-1998:Past}, grey line and shading show total solar insolation estimate with 1$\sigma$ uncertainty \cite[]{Steinhilber-Beer-2009:Total}, and red lines show global stratospheric sulfate aerosol injection \cite[]{gao2008}.

This figure is definitely too busy, I just thought it would be interesting to compare these records to some common forcing plot and see if anything pops out. 

[For discussion: ``old'' solution reconstructs 0.9 $\pm$ 2.7K Nordic Sea cooling from 800-1860CE, and 0.5 $\pm$ 2.2K subpolar gyre cooling from 1230-1900CE. ]}
\end{figure}




\section{Conclusion}

%
% The main text should start with an introduction. Except for short
% manuscripts (such as comments and replies), the text should be divided
% into sections, each with its own heading.

% Headings should be sentence fragments and do not begin with a
% lowercase letter or number. Examples of good headings are:

% \section{Materials and Methods}
% Here is text on Materials and Methods.
%
% \subsection{A descriptive heading about methods}
% More about Methods.
%
% \section{Data} (Or section title might be a descriptive heading about data)
%
% \section{Results} (Or section title might be a descriptive heading about the
% results)
%
% \section{Conclusions}



%Text here ===>>>


%%

%  Numbered lines in equations:
%  To add line numbers to lines in equations,
%  \begin{linenomath*}
%  \begin{equation}
%  \end{equation}
%  \end{linenomath*}



%% Enter Figures and Tables near as possible to where they are first mentioned:
%
% DO NOT USE \psfrag or \subfigure commands.
%
% Figure captions go below the figure.
% Acronyms used in figure captions will be spelled out in the final, published version.

% Table titles go above tables;  other caption information
%  should be placed in last line of the table, using
% \multicolumn2l{$^a$ This is a table note.}
% NOTE that there is no difference between table caption and table heading in the final, published version
%
%----------------
% EXAMPLE FIGURES
%
% \begin{figure}
% \includegraphics{example.png}
% \caption{caption}
% \end{figure}
%
% Giving latex a width will help it to scale the figure properly. A simple trick is to use \textwidth. Try this if large figures run off the side of the page.
% \begin{figure}
% \noindent\includegraphics[width=\textwidth]{anothersample.png}
%\caption{caption}
%\label{pngfiguresample}
%\end{figure}
%
%
% If you get an error about an unknown bounding box, try specifying the width and height of the figure with the natwidth and natheight options. This is common when trying to add a PDF figure without pdflatex.
% \begin{figure}
% \noindent\includegraphics[natwidth=800px,natheight=600px]{samplefigure.pdf}
%\caption{caption}
%\label{pdffiguresample}
%\end{figure}
%
%
% PDFLatex does not seem to be able to process EPS figures. You may want to try the epstopdf package.
%

%
% ---------------
% EXAMPLE TABLE
%
% \begin{table}
% \caption{Time of the Transition Between Phase 1 and Phase 2$^{a}$}
% \centering
% \begin{tabular}{l c}
% \hline
%  Run  & Time (min)  \\
% \hline
%   $l1$  & 260   \\
%   $l2$  & 300   \\
%   $l3$  & 340   \\
%   $h1$  & 270   \\
%   $h2$  & 250   \\
%   $h3$  & 380   \\
%   $r1$  & 370   \\
%   $r2$  & 390   \\
% \hline
% \multicolumn{2}{l}{$^{a}$Footnote text here.}
% \end{tabular}
% \end{table}

%%%%%%%%%%%%%%%%%%%%%%%%%%%%%%%%%%%%%%%%%%%%%%%
% SIDEWAYS FIGURES and TABLES
% AGU prefers the use of {sidewaystable} over {landscapetable} as it causes fewer problems.
%
% \begin{sidewaysfigure}
% \includegraphics[width=20pc]{figsamp}
% \caption{caption here}
% \label{newfig}
% \end{sidewaysfigure}
%
%  \begin{sidewaystable}
%  \caption{Caption here}
% \label{tab:signif_gap_clos}
%  \begin{tabular}{ccc}
% one&two&three\\
% four&five&six
%  \end{tabular}
%  \end{sidewaystable}

%% If using numbered lines, please surround equations with \begin{linenomath*}...\end{linenomath*}
%\begin{linenomath*}
%\begin{equation}
%y|{f} \sim g(m, \sigma),
%\end{equation}
%\end{linenomath*}

%%% End of body of article

%%%%%%%%%%%%%%%%%%%%%%%%%%%%%%%%%%%%%%%%%%%%%%%
%% Optional Appendices go here
%
% The \appendix command resets counters and redefines section heads
%
% After typing \appendix
%
%\section{Here Is Appendix Title}
% will show
% A: Here Is Appendix Title
%
%\appendix
%\section{Here is a sample appendix}

%%%%%%%%%%%%%%%%%%%%%%%%%%%%%%%%%%%%%%%%%%%%%%%
% Optional Glossary, Notation or Acronym section goes here:
%
% Glossary is only allowed in Reviews of Geophysics
%  \begin{glossary}
%  \term{Term}
%   Term Definition here
%  \term{Term}
%   Term Definition here
%  \term{Term}
%   Term Definition here
%  \end{glossary}


%%%%%%%%%%%%%%%%%%%%%%%%%%%%%%%%%%%%%%%%%%%%%%%
% Acronyms
%% NOTE that acronyms in the final published version will be spelled out when used in figure captions.
%   \begin{acronyms}
%   \acro{Acronym}
%   Definition here
%   \acro{EMOS}
%   Ensemble model output statistics
%   \acro{ECMWF}
%   Centre for Medium-Range Weather Forecasts
%   \end{acronyms}


%%%%%%%%%%%%%%%%%%%%%%%%%%%%%%%%%%%%%%%%%%%%%%%
% Notation
%   \begin{notation}
%   \notation{$a+b$} Notation Definition here
%   \notation{$e=mc^2$}
%   Equation in German-born physicist Albert Einstein's theory of special
%  relativity that showed that the increased relativistic mass ($m$) of a
%  body comes from the energy of motion of the body—that is, its kinetic
%  energy ($E$)—divided by the speed of light squared ($c^2$).
%   \end{notation}




%%%%%%%%%%%%%%%%%%%%%%%%%%%%%%%%%%%%%%%%%%%%%%%
%
% DATA SECTION and ACKNOWLEDGMENTS
%
%%%%%%%%%%%%%%%%%%%%%%%%%%%%%%%%%%%%%%%%%%%%%%%

\section*{Open Research Section}
This section MUST contain a statement that describes where the data supporting the conclusions can be obtained. Data cannot be listed as ''Available from authors'' or stored solely in supporting information. Citations to archived data should be included in your reference list. Wiley will publish it as a separate section on the paper’s page. Examples and complete information are here:
https://www.agu.org/Publish with AGU/Publish/Author Resources/Data for Authors


\acknowledgments
Enter acknowledgments here. This section is to acknowledge funding, thank colleagues, enter any secondary affiliations, and so on.

\clearpage
%%%%%%%%%%%%%%%%%%%%%%%%%%%%%%%%%%%%%%%%%%%%%%%
% REFERENCES and BIBLIOGRAPHY
%
\bibliography{Hamilton2024} 
% don't specify bibliographystyle
%
%%%%%%%%%%%%%%%%%%%%%%%%%%%%%%%%%%%%%%%%%%%%%%%


%Reference citation instructions and examples:
%
% Please use ONLY \cite and \citeA for reference citations.
% \cite for parenthetical references
% ...as shown in recent studies (Simpson et al., 2019)
% \citeA for in-text citations
% ...Simpson et al. (2019) have shown...
%
%
%...as shown by \citeA{jskilby}.
%...as shown by \citeA{lewin76}, \citeA{carson86}, \citeA{bartoldy02}, and \citeA{rinaldi03}.
%...has been shown \cite{jskilbye}.
%...has been shown \cite{lewin76,carson86,bartoldy02,rinaldi03}.
%... \cite <i.e.>[]{lewin76,carson86,bartoldy02,rinaldi03}.
%...has been shown by \cite <e.g.,>[and others]{lewin76}.
%
% apacite uses < > for prenotes and [ ] for postnotes
% DO NOT use other cite commands (e.g., \citet, \citep, \citeyear, \nocite, \citealp, etc.).
%


\clearpage
\end{document}



More Information and Advice:

%%%%%%%%%%%%%%%%%%%%%%%%%%%%%%%%%%%%%%%%%%%%%%%
%
%  SECTION HEADS
%
%%%%%%%%%%%%%%%%%%%%%%%%%%%%%%%%%%%%%%%%%%%%%%%

% Capitalize the first letter of each word (except for
% prepositions, conjunctions, and articles that are
% three or fewer letters).

% AGU follows standard outline style; therefore, there cannot be a section 1 without
% a section 2, or a section 2.3.1 without a section 2.3.2.
% Please make sure your section numbers are balanced.
% ---------------
% Level 1 head
%
% Use the \section{} command to identify level 1 heads;
% type the appropriate head wording between the curly
% brackets, as shown below.
%
%An example:
%\section{Level 1 Head: Introduction}
%
% ---------------
% Level 2 head
%
% Use the \subsection{} command to identify level 2 heads.
%An example:
%\subsection{Level 2 Head}
%
% ---------------
% Level 3 head
%
% Use the \subsubsection{} command to identify level 3 heads
%An example:
%\subsubsection{Level 3 Head}
%
%---------------
% Level 4 head
%
% Use the \subsubsubsection{} command to identify level 3 heads
% An example:
%\subsubsubsection{Level 4 Head} An example.
%
%%%%%%%%%%%%%%%%%%%%%%%%%%%%%%%%%%%%%%%%%%%%%%%
%
%  IN-TEXT LISTS
%
%%%%%%%%%%%%%%%%%%%%%%%%%%%%%%%%%%%%%%%%%%%%%%%
%
% Do not use bulleted lists; enumerated lists are okay.
% \begin{enumerate}
% \item
% \item
% \item
% \end{enumerate}
%
%%%%%%%%%%%%%%%%%%%%%%%%%%%%%%%%%%%%%%%%%%%%%%%
%
%  EQUATIONS
%
%%%%%%%%%%%%%%%%%%%%%%%%%%%%%%%%%%%%%%%%%%%%%%%

% Single-line equations are centered.
% Equation arrays will appear left-aligned.

Math coded inside display math mode \[ ...\]
 will not be numbered, e.g.,:
 \[ x^2=y^2 + z^2\]

 Math coded inside \begin{equation} and \end{equation} will
 be automatically numbered, e.g.,:
 \begin{equation}
 x^2=y^2 + z^2
 \end{equation}


% To create multiline equations, use the
% \begin{eqnarray} and \end{eqnarray} environment
% as demonstrated below.
\begin{eqnarray}
  x_{1} & = & (x - x_{0}) \cos \Theta \nonumber \\
        && + (y - y_{0}) \sin \Theta  \nonumber \\
  y_{1} & = & -(x - x_{0}) \sin \Theta \nonumber \\
        && + (y - y_{0}) \cos \Theta.
\end{eqnarray}

%If you don't want an equation number, use the star form:
%\begin{eqnarray*}...\end{eqnarray*}

% Break each line at a sign of operation
% (+, -, etc.) if possible, with the sign of operation
% on the new line.

% Indent second and subsequent lines to align with
% the first character following the equal sign on the
% first line.

% Use an \hspace{} command to insert horizontal space
% into your equation if necessary. Place an appropriate
% unit of measure between the curly braces, e.g.
% \hspace{1in}; you may have to experiment to achieve
% the correct amount of space.


%%%%%%%%%%%%%%%%%%%%%%%%%%%%%%%%%%%%%%%%%%%%%%%
%
%  EQUATION NUMBERING: COUNTER
%
%%%%%%%%%%%%%%%%%%%%%%%%%%%%%%%%%%%%%%%%%%%%%%%

% You may change equation numbering by resetting
% the equation counter or by explicitly numbering
% an equation.

% To explicitly number an equation, type \eqnum{}
% (with the desired number between the brackets)
% after the \begin{equation} or \begin{eqnarray}
% command.  The \eqnum{} command will affect only
% the equation it appears with; LaTeX will number
% any equations appearing later in the manuscript
% according to the equation counter.
%

% If you have a multiline equation that needs only
% one equation number, use a \nonumber command in
% front of the double backslashes (\\) as shown in
% the multiline equation above.

% If you are using line numbers, remember to surround
% equations with \begin{linenomath*}...\end{linenomath*}

%  To add line numbers to lines in equations:
%  \begin{linenomath*}
%  \begin{equation}
%  \end{equation}
%  \end{linenomath*}




%%% Local Variables:
%%% mode: latex
%%% TeX-master: t
%%% End:

\message{ !name(Hamilton2024.tex) !offset(-829) }
